\section{Midpoint transfer rules}\label{met:transfer-section}
The principal forest bound used an arbitrary-displacement quadratic
estimate. We now give transfers for other intrinsic inputs, to be used
in the path-cost and recursive models below. The three mechanisms are
additive loss of a power of the slope, exclusion of separated lenses,
and finite-codimensional gain. The approximation head is fixed before
selecting a pair of deviations; no child supremum is assumed attained.

\subsection{Homogeneous tail estimates and separated lenses}
\begin{proposition}[Homogeneous tail estimates]\label{met:tail-to-lens}
Let $X$ have a specified coordinate system with dense finite-support
vectors. Fix a family $\mathcal P$ of bounded finite-rank coordinate
projections such that every finite support is contained in the range of
some $P\in\mathcal P$. Consider separately the following two hypotheses,
each imposed for every $P\in\mathcal P$, all pure head and tail vectors
$p=Pp$, $Py=0$, and every $R>0$:
\begin{align}
 \|y\|^2&\le C^2 R(R-\|p\|)
 &&\text{if }\|p\pm y\|\le R, \label{met:radial-input}\\
 \|p\|^2+a\|y\|^2&\le R^2
 &&\text{if }\|p\pm y\|\le R. \label{met:square-input}
\end{align}
Here $C>0$, respectively $a>0$, is uniform over $\mathcal P$, the
vectors, and $R$. No uniform bound on $\|P\|$ or $\|I-P\|$ is assumed.
For $R>0$ and $\varepsilon>0$, an infinite $\varepsilon$-separated family in the symmetric $R$-lens at
$x$ then satisfies, respectively,
\begin{align}
 \varepsilon^2/4&\le C^2R(R-\|x\|),\label{met:radial-lens}\\
 \|x\|^2+a\varepsilon^2/4&\le R^2.\label{met:square-lens}
\end{align}
\end{proposition}
\begin{proof}
Choose $p$ and $P$ as in Lemma~\ref{met:finite-pair}, and apply the
assumed estimate to $p,y$ at radius $R+2\eta$. In the first case,
\[
 (\varepsilon/2-\eta)^2
 \le C^2(R+2\eta)(R+2\eta-\|p\|)
 \le C^2(R+2\eta)(R+3\eta-\|x\|).
\]
In the second case,
$\max\{\|x\|-\eta,0\}^2+a(\varepsilon/2-\eta)^2\le(R+2\eta)^2$.
Take $0<\eta<\varepsilon/2$ and let $\eta$ decrease to zero.
\end{proof}

\subsection{Squared losses and additive losses}\label{met:squared-section}
For the noncontracting normalization define
\[
 R_j=\sup\{\|x_e\|: e\text{ is a stage-$j$ edge}\},
 \qquad 1\le R_j\le D.
\]
The countable supremum need not be attained. Our estimates apply to each
parent edge before the supremum is taken.

The squared-lens hypothesis below is the $p=2$ Banach-space form of
the metric diamond inequality in
\cite[Equation~(2) and Theorem~3.7 of arXiv v2]{Malthaner2025}:
its endpoints are $-x,x$ and its intermediate points are the deviations
$z_i$. We retain the coefficient by summing the squared losses directly.

\begin{theorem}\label{met:squared-transfer}
If every infinite $\varepsilon$-separated symmetric $R$-lens satisfies
$\|x\|^2\le R^2-b\varepsilon^2$ for some $b>0$, then every embedding
\eqref{met:all-pairs} satisfies $D^2\ge1+4bk$.
If instead \eqref{met:radial-lens} holds, then
\begin{equation}\label{met:additive-bound}
 k\le C^2D(D-1).
\end{equation}
\end{theorem}
\begin{proof}
For each stage-$j$ parent, its children lie in the lens of radius
$R_{j+1}$ and its deviations are $2$-separated. The first hypothesis
gives $\|x\|^2\le R_{j+1}^2-4b$. Taking suprema gives
$R_j^2\le R_{j+1}^2-4b$, and summing yields
$1\le R_0^2\le D^2-4bk$.
For the radial estimate,
\[
 1\le C^2R_{j+1}(R_{j+1}-\|x\|),\qquad
 R_j\le R_{j+1}-\frac1{C^2R_{j+1}}
 \le R_{j+1}-\frac1{C^2D}.
\]
Summing proves \eqref{met:additive-bound}. Empty sums cover $k=0$.
\end{proof}

\subsection{Retaining a homogeneous power loss}
The radius dependence of a higher-power input can be accumulated directly.
Replacing it by a fixed contraction factor would discard this information.

\begin{theorem}\label{met:power-transfer}
Assume the density and projection-family hypotheses of
Proposition~\ref{met:tail-to-lens}. Let $q\ge2$ and $c>0$. Suppose that
for every admissible projection, every pure head--tail pair $p,y$, and
every $R>0$,
\[
 \|p\pm y\|\le R
 \quad\Longrightarrow\quad
 \|p\|^2+cR^{2-q}\|y\|^q\le R^2.
\]
Then every distortion-$D$ embedding of $\mathcal D_k$ satisfies
\[
 ck\le D^{q-2}(D^2-1).
\]
\end{theorem}
\begin{proof}
Use the noncontracting normalization. Fix a parent $x$ whose child
slopes have norm at most $R$. Lemma~\ref{met:finite-pair}, followed
by $\eta\downarrow0$, gives
\[
 \|x\|^2+cR^{2-q}\le R^2.
\]
Here the paired tail has limiting lower norm one. Set $R=R_{j+1}$
and take the supremum over stage-$j$ parents. Since $1\le R_{j+1}\le D$
and $q\ge2$,
\[
 R_{j+1}^2-R_j^2\ge cR_{j+1}^{2-q}\ge cD^{2-q}.
\]
Summing between $R_0\ge1$ and $R_k\le D$ gives
$D^2-1\ge kcD^{2-q}$, as required.
\end{proof}

\subsection{Keeping the parameterized test losses}\label{met:parametric-section}
Some primal tests give a linear gain before optimization. The graph
argument can preserve that parameter instead of replacing the input by a
single modulus.

\begin{proposition}\label{met:parametric}
Assume the density and projection-family hypotheses of
Proposition~\ref{met:tail-to-lens}. Fix constants $a,b>0$. Suppose that
for every admissible projection and every pure head--tail pair $p,y$,
the following estimate holds for all $\lambda>0$:
\[
 \max_\pm\|p\pm y\|\ge
 \frac{\|p\|+a\lambda\|y\|}{\sqrt{1+b\lambda^2}}.
\]
Put $Q=\sqrt{1+b\lambda^2}$. For noncontracting diamond slopes
bounded above by $D$, the stage suprema satisfy
\[
 R_j\le Q R_{j+1}-a\lambda
 \le R_{j+1}-\bigl(a\lambda-(Q-1)D\bigr).
\]
In particular, whenever $\Delta=a\lambda-(Q-1)D>0$, one has
$1+k\Delta\le D$.
\end{proposition}
\begin{proof}
For each parent $x$, choose a finite-support approximation $p$ and a
paired tail $y$ as in Lemma~\ref{met:finite-pair}. The half-difference
has norm at least one, so $\|y\|\ge1-\eta$. If the child bound is
$R$, then $R+2\eta\ge(\|p\|+a\lambda(1-\eta))/Q$.
Let $\eta\downarrow0$ to obtain $\|x\|\le QR-a\lambda$,
and then take suprema. Summation proves the last assertion.
\end{proof}

\subsection{Backward contraction}\label{met:backward-section}
An exact power-law lens estimate is useful but unnecessary for a
qualitative obstruction. The following form also handles estimates whose
constants are selected only after the separation is specified.

\begin{proposition}\label{met:backward}
Fix $0<\alpha\le1$. Suppose some $0\le\rho<1$ has the following
property: whenever $\|x\pm z_i\|\le1$ and
$\|z_i-z_j\|\ge2\alpha$ for $i\ne j$, one has $\|x\|\le\rho$.
An embedding with lower factor $\alpha$ and upper factor one must satisfy
\begin{equation}\label{met:backward-bound}
 \alpha\le\rho^k,
\end{equation}
where the empty product $\rho^0$ is one.
\end{proposition}
\begin{proof}
Let $r_j$ be the supremum of stage-$j$ slope norms. Initially $r_k\le1$,
and every $r_j\ge\alpha>0$. If all child slopes are at most $r_{j+1}$,
normalize \eqref{met:slopes} by this positive number. The resulting
unit lens has separation at least $2\alpha/r_{j+1}\ge2\alpha$.
The hypothesis gives $\|x\|\le\rho r_{j+1}$ for every parent, hence
$r_j\le\rho r_{j+1}$. Induction proves the assertion. If $\rho=0$,
the existence of any replacement already contradicts $r_j\ge\alpha$.
\end{proof}

For example, \eqref{met:radial-lens} gives
$\rho=1-\alpha^2/C^2$ when this number is positive. If it is
nonpositive, a parent lens with center norm at least $\alpha>0$
is already impossible. With upper factor one it also gives the stronger
intermediate inequality
\[
 \|x\|\le r-\frac{\alpha^2}{C^2r}
 \le r-\frac{\alpha^2}{C^2},\qquad 0<r\le1.
\]
Thus $k\le C^2(D^2-D)$ when $\alpha=1/D$. Alternatively, if the proved
lens input is $\|x\|\le1-\eta(2\alpha)$, the direct conclusion is
$\alpha\le(1-\eta(2\alpha))^k$; a polynomial rate requires an actual
quantitative bound for $\eta$.

\subsection{Enlarged lenses and stable parents}\label{met:selection-section}
\label{met:bridge:enlarged-section}
A unit-ball lens gap bounds the parent relative to the child bound.
The equivalent normalization below divides by the parent norm instead.
It also identifies an actual parent at which all children have moderate
stretch.

\begin{proposition}[Enlarged lenses and a stable parent]
\label{met:stable-parent}\label{met:bridge:enlarged-backward}
Let $D\ge1$ and $\delta>0$. If $(1+\delta)^k>D$, every
noncontracting embedding with upper factor $D$ has a parent edge all of whose
child slopes have norm at most $(1+\delta)$ times its own slope norm.
If, for every unit $u$, the enlarged lens
\[
 M(u,\delta)=\{z:\max_\pm\|u\pm z\|\le1+\delta\}
\]
contains no infinite $2/D$-separated family, then every such embedding
satisfies
\begin{equation}\label{met:bridge:enlarged-depth}
 1\le \frac{D}{(1+\delta)^k}.
\end{equation}
\end{proposition}
\begin{proof}
If no stable parent existed, choose at each stage an actual child whose
slope norm is greater than $1+\delta$ times its parent's norm. Starting
with norm at least one would contradict the terminal upper bound $D$.
At a stable parent of norm $r\le D$, division by $r$ puts all deviations
in $M(x/r,\delta)$ with separation at least $2/r\ge2/D$.
This proves the second assertion by contradiction.

The same exclusion also gives a backward proof. If all children of a
parent of norm $r$ have norm at most $L\le D$, then
$r\ge L/(1+\delta)$ would put the normalized deviations in the forbidden
lens. Thus every parent has norm less than $L/(1+\delta)$; taking
suprema gives the non-strict bound, and iteration from $L=D$ proves
\eqref{met:bridge:enlarged-depth}. Neither proof assumes attainment of a
child supremum.
\end{proof}

The stable parent may instead be divided by its child bound
$(1+\delta)r$. Its unit-ball lens has center norm $1/(1+\delta)$
and separation at least $2/((1+\delta)D)$. This is useful when the
intrinsic estimate is stated for unit-ball lenses rather than enlarged
unit-center lenses.

\subsection{Uniform finite-head bounds and their parameters}\label{met:packing-section}
Many estimates above control all lens vectors using one head selected
from the center. This order of quantifiers is precisely what packing
requires.

\begin{lemma}\label{met:packing}
Let $P$ be bounded and finite rank, let $A\subset X$ be bounded,
and let $b\ge0$ and $\varepsilon>0$.
If $\|(I-P)v\|\le b$ for all $v\in A$, then $A$ contains no infinite
$\varepsilon$-separated family when $2b<\varepsilon$.
Consequently the following two inputs have different graph conclusions.
\begin{enumerate}
\item Let $0<\delta<1$. For each $\|x\|>1-\delta$, suppose one such $P$ bounds by
$b(\delta)$ the tails of every $v$ with $\|x\pm v\|\le1$.
If $2b(\delta)<2/D$, backward contraction gives
$D^{-1}\le(1-\delta)^k$.
\item Let $\delta>0$. For every unit $x$, suppose one such $P$ bounds by
$b(\delta)$ the tails of every $v$ with $\|x\pm v\|\le1+\delta$.
If $2b(\delta)<2/D$, actual-child selection gives
$D\ge(1+\delta)^k$.
\end{enumerate}
The same conclusions hold when the bound is on the first endpoint of
every pair with $(u+v)/2=x$: take $u,v\in B_X$ in case (1), and
$u,v\in(1+\delta)B_X$ in case (2).
\end{lemma}
\begin{proof}
If $(v_i)$ were such a separated family, choose distinct $i,j$ with
$\|Pv_i-Pv_j\|<\varepsilon-2b$. The triangle inequality contradicts
$\|v_i-v_j\|\ge\varepsilon$. The graph conclusions follow from
Propositions~\ref{met:backward} and~\ref{met:stable-parent}. For first
endpoints, translation by their common center preserves separation and
the same finite-dimensional argument applies directly to $(u_i)$.
\end{proof}

Under the projection-family hypotheses of Proposition~\ref{met:tail-to-lens},
a bound only on pure tails uses Lemma~\ref{met:finite-pair} for a
second form of compactness. For example, suppose unit-ball
head--tail pairs obey $\|y\|\le g(1-\|p\|)$, where
$g:[0,1]\to[0,\infty)$ is nondecreasing and $g(r)\to0$ as
$r\downarrow0$. For $t>0$, an infinite
$2t$-separated unit lens cannot have center norm greater than $1-\Delta$
if
\begin{equation}\label{met:g-parameters}
 0<\Delta<1/4,\qquad (1+2\Delta)g(4\Delta)+\Delta<t.
\end{equation}
Indeed choose a head approximating the center within $\Delta$ and a
half-difference with head norm at most $\Delta$. After division by
$1+2\Delta$, its tail norm is at least $(t-\Delta)/(1+2\Delta)$,
whereas the normalized center deficit is at most $4\Delta$. This
contradicts~\eqref{met:g-parameters}. Thus with $t=1/D$ the resulting
backward factor is $1-\Delta$. Neither monotonicity nor a power-law rate
is presumed for a function until its actual formula establishes it.

\subsection{From a qualitative sequential estimate to a uniform lens gap}
\label{met:bridge:sequential-section}
The qualitative argument can start with a statement about sequences of
pure head--tail pairs, without first constructing a numerical modulus.


\begin{lemma}\label{met:bridge:sequential-lens}
Assume the density and projection-family hypotheses of
Proposition~\ref{met:tail-to-lens}. Suppose that whenever admissible heads $P_n\in\mathcal P$
and vectors $p_n,y_n$ satisfy
\[
 P_np_n=p_n,\qquad P_ny_n=0,\qquad
 \|p_n\pm y_n\|\le1,\qquad \|p_n\|\longrightarrow1,
\]
one has $\|y_n\|\longrightarrow0$. The heads are allowed to vary
with $n$. Then, for each $\varepsilon>0$, there is $0<\rho<1$
such that every infinite $\varepsilon$-separated symmetric unit lens
has center norm at most $\rho$.
\end{lemma}
\begin{proof}
If no such $\rho$ existed, choose centers $x_n$ with
$\|x_n\|\to1$ and infinite $\varepsilon$-separated families
$(z_{n,i})_i$ with $\|x_n\pm z_{n,i}\|\le1$.
Choose $\eta_n\downarrow0$, finitely supported $p_n$ with
$\|p_n-x_n\|<\eta_n$, and admissible heads $P_n$ satisfying
$P_np_n=p_n$. Lemma~\ref{met:finite-pair} gives a half-difference
$w_n$ and the tail $y_n=w_n-P_nw_n$ with
\[
 \|y_n\|\ge\varepsilon/2-\eta_n,\qquad
 \|p_n\|\ge\|x_n\|-\eta_n,\qquad
 \|p_n\pm y_n\|\le1+2\eta_n.
\]
Divide both vectors by $1+2\eta_n$. The resulting pure head--tail
pairs have endpoints in the unit ball and centers whose norms tend
to one. Their tail norms have lower limit at least
$\varepsilon/2>0$, contradicting the assumed sequential statement.
This contradiction gives one $\rho$ valid for all centers and heads.
\end{proof}

In the upper-one graph normalization, take $\varepsilon=2/D$.
Dividing a parent configuration by any child bound $0<L\le1$
leaves separation at least $2/(DL)\ge\varepsilon$. Thus the lemma
gives the parent bound $\rho L$ and
Proposition~\ref{met:backward} yields $1/D\le\rho^k$.
\subsection{Finite-codimensional gains}\label{met:quotient-section}
The relation between finite-codimensional midpoint gains and small
symmetric lenses is developed in \cite[Theorem~2.1]{DKRRZ2016}.
Here we retain the quantitative difference between selecting an actual
child with a reserved error margin and taking a limiting supremum.
The next transfer uses the symmetric maximum, so it also applies when the
intrinsic theorem controls the average of the endpoint norms.

\begin{proposition}\label{met:quotient-transfer}
Suppose that for some $t>0$ and $g>0$, every unit vector $u\in X$ admits
a closed finite-codimensional subspace $F_u$ such that
\[
 y\in F_u,\quad\|y\|\ge t
 \quad\Longrightarrow\quad
 \max\{\|u+y\|,\|u-y\|\}\ge1+g.
\]
If $D\ge1$ and $t\le1/(2D)$, every noncontracting embedding of
$\mathcal D_k$ with upper factor $D$ satisfies $D\ge(1+g/2)^k$.
\end{proposition}
\begin{proof}
For a parent slope $x$, put $r=\|x\|\in[1,D]$ and use the subspace
$F_{x/r}$. Apply the quotient part of Lemma~\ref{met:finite-pair} to
choose $w=(z_i-z_j)/2$, with $\|w\|\ge1$, and $w'\in F_{x/r}$
such that $\|w-w'\|<\min\{1/2,rg/2\}$. Then
$\|w'/r\|\ge1/(2D)\ge t$, and
\[
 \max_\pm\|x\pm w\|\ge r(1+g)-\|w-w'\|>r(1+g/2).
\]
Each $x\pm w$ is the average of two of the four selected child slopes.
One of these four actual slopes is therefore greater than $r(1+g/2)$.
Choose it and repeat through $k$ generations. This argument makes no
claim that a supremum over all children is attained.
\end{proof}

A bound initially stated only for $\|y\|=t$ extends to $\|y\|\ge t$:
for a fixed direction $v$, the function
$h(r)=\max\{\|u+rv\|,\|u-rv\|\}$ is even and convex, and hence
nondecreasing on $[0,\infty)$. This observation applies likewise to the
average of the two norms.

\begin{proposition}[The full asymptotic gain in a backward step]
\label{met:full-gain}
Assume the finite-codimensional hypothesis of
Proposition~\ref{met:quotient-transfer}. Consider a noncontracting embedding
with upper factor $D\ge1$. If $t\le1/D$ and all child
slopes of one parent have norm at most $L$, then
$\|x\|\le L/(1+g)$. Consequently every depth-$k$ embedding with
upper factor $D$ satisfies $1\le D/(1+g)^k$.
\end{proposition}
\begin{proof}
Put $r=\|x\|\in[1,D]$. Quotient compactness supplies half-differences
$w$ with $\|w\|\ge1$ and vectors $v\in F_{x/r}$ with
$\|v-w\|<\eta$, where $0<\eta<1/2$ is arbitrarily small.
Set $w'=\|w\|v/\|v\|$. Then $w'\in F_{x/r}$,
$\|w'\|/r=\|w\|/r\ge1/D\ge t$, and
\[
 \|w'-w\|\le \bigl|\|w\|-\|v\|\bigr|+\|v-w\|<2\eta.
\]
The intrinsic gain and the crossed-average child bound imply
$(1+g)r\le\max_\pm\|x\pm w'\|\le L+2\eta$.
Let $\eta\downarrow0$ before taking suprema over parents, then
iterate backward. Without the norm-preserving rescaling, the same
proof works under the smaller threshold $t\le1/(2D)$.
\end{proof}

\begin{lemma}[The local midpoint estimate gives an enlarged lens gap]
\label{met:bridge:local-amuc}
Let $R>0$ and put $t=R/4$. Suppose $\gamma>0$ is such that every
unit $u$ admits a closed finite-codimensional subspace $F_u$ with
\begin{equation}\label{met:bridge:local-amuc-input}
 \max_\pm\|u\pm t y\|\ge1+\gamma
 \qquad(y\in F_u,\ \|y\|=1).
\end{equation}
Then $\beta=\gamma/2$ excludes infinite $R$-separated families in
$M(u,\beta)$, uniformly over unit $u$.
\end{lemma}
\begin{proof}
If $(z_i)\subset M(u,\beta)$ were such a family, symmetry and
convexity would place every half-difference $(z_i-z_j)/2$ in the
same lens. The family is bounded, so its image in the
finite-dimensional quotient $X/F_u$ has arbitrarily close distinct
pairs. Hence there are $i\ne j$ and $w\in F_u$ such that, with
$v=(z_i-z_j)/2$,
\[
 \|v\|\ge R/2,\qquad
 \|v-w\|<\min\{R/4,\gamma/2\}.
\]
In particular $\|w\|>R/4=t$. For a fixed unit direction $y$, the
function $s\mapsto\max_\pm\|u\pm s y\|$ is even and convex,
hence nondecreasing for $s\ge0$. Apply
\eqref{met:bridge:local-amuc-input} to $y=w/\|w\|$ to obtain
$\max_\pm\|u\pm w\|\ge1+\gamma$. On the other hand,
$v\in M(u,\beta)$ and the strict approximation error give
\[
 \max_\pm\|u\pm w\|
 \le1+\beta+\|w-v\|<1+\gamma,
\]
a contradiction. For $R=2/D$, Proposition~\ref{met:bridge:enlarged-backward}
now gives \eqref{met:bridge:enlarged-depth} with $\delta=\beta=\gamma/2$.
\end{proof}

\subsection{The obstruction for an equivalent midpoint-convex norm}
\label{met:bridge:renorming-section}
The following obstruction is the countably branching diamond case of
\cite[Theorem~4.1]{BaudierEtAl2017}. We include the short deduction
from the quotient transfer to specify the norm convention and the
effect of an equivalent norm.

\begin{theorem}\label{met:bridge:equivalent-amuc}
If a Banach space admits an equivalent asymptotically midpoint
uniformly convex norm, then the family $(\mathcal D_k)_{k\ge0}$
does not embed into it with uniformly bounded distortion.
Write $\operatorname{cof}(X)$ for the closed finite-codimensional
linear subspaces of $X$. In infinite dimension, midpoint convexity may be expressed by
positivity of the maximum modulus
\begin{equation}\label{met:bridge:max-modulus}
 \delta_N^{\max}(t)=
 \inf_{N(u)=1}\ 
 \sup_{F\in\operatorname{cof}(X)}\ 
 \inf_{\substack{y\in F\\N(y)=1}}
 \bigl(\max_\pm N(u\pm t y)-1\bigr)>0
 \qquad(0<t<1).
\end{equation}
\end{theorem}
\begin{proof}
In finite dimension, even $\mathcal D_1$ has a bounded infinite
separated family of midpoints, whose images under a bi-Lipschitz
map would contradict total boundedness. Consider an
infinite-dimensional space, and suppose $a\|x\|\le N(x)\le b\|x\|$
with $a,b>0$. Embeddings with common distortion $D_0$ for the original
norm, and arbitrary positive scales, would give embeddings with common
distortion $D=(b/a)D_0$ for $N$: divide each map by its original
lower scale times $a$.

Fix $t=1/(2D)$ and $0<g<\delta_N^{\max}(t)$. The definition of the
supremum in \eqref{met:bridge:max-modulus} supplies, for every
$N$-unit $u$, a finite-codimensional $F_u$ on whose unit sphere the
maximum is at least $1+g$. Even convexity in each direction extends
this bound to every vector of $F_u$ with $N$-norm at least $t$.
Proposition~\ref{met:quotient-transfer}, applied to $N$, therefore
gives $D\ge(1+g/2)^k$ for every $k$, which is impossible.
An averaged midpoint lower bound also supplies the required maximum
bound, since the maximum dominates the average.
\end{proof}


More generally, if $a\|v\|\le N(v)\le b\|v\|$, an embedding with
distortion $D$ for $N$ has distortion at most $(b/a)D$ for the original
norm, and conversely with the same factor. Thus the absence of uniform
diamond distortion is invariant under equivalent renorming, while the
numerical constants attached to a given norm need not be.
